\documentclass{article}
\usepackage[utf8]{inputenc}

\usepackage{graphicx}
\usepackage[margin=0.5in]{geometry}
\usepackage{array}
\usepackage{multirow}
\usepackage{hyperref}
\usepackage{xcolor}

\title{DSAA-6000Q Project Final Report Template}
\author{Team Members: [1-5 people] \\ Date: [Submission Date]}
\date{}

\usepackage{natbib}
\usepackage{graphicx}

\begin{document}

\maketitle



\begin{abstract}
A brief summary of the project, including the research question or objective, methods, and key findings.
\end{abstract}
\line(1,0){550} 
\section*{Timeline [Important!]}
\begin{itemize}
    \item \textbf{Team-up due}: {\textbf{\color{red} April 19th, 2026.}} 
    
    Please complete the team-up and register at this link \href{https://docs.google.com/spreadsheets/d/1FJsGVfhf7oHint2zKoZ13pNIbClVsMkPLsbs6yYbZAE/edit?usp=sharing}{[Google Sheet (Teams)]}.
    \item \textbf{Project topic \& TA-matching due}: {\textbf{\color{red} April 19th, 2026.}}

    Please decide on the topic of your project, fill in the column of ``Project name" in this link \href{https://docs.google.com/spreadsheets/d/1FJsGVfhf7oHint2zKoZ13pNIbClVsMkPLsbs6yYbZAE/edit?usp=sharing}{[Google Sheet (Teams)]}.

    Please also reach out to 1-2 TAs about your project plan via e-mail, budget requirements of the potential resources required in your project (if any). 
    
    The TA would be happy to help with supervising / question-answering your project.  

    {\textbf{\color{red} Self-proposed project requires approval of the instructor.}}
    \item \textbf{Project presentation due}: {\textbf{\color{red} May 12th, 2026.}}

    In the week-13 (final 2 lectures (May 11, May 12)), each team will have around 8 minutes to present and share their project results.

    The presentation grade will consist of 20\% of your project grades.
    \item \textbf{Final report \& supplementation materials due}: {\textbf{\color{red} May 17th, 2026.}}

    Last day to submit the final project report of your team, {\textbf{\color{red} no late days are allowed.}}
\end{itemize}

\line(1,0){550} 
\section{Introduction}
\begin{itemize}
    \item Background and context of the project.
    \item Motivation and significance of the study.
    \item Objectives and research questions.
\end{itemize}

\section{Related Works}
\begin{itemize}
    \item Summary of existing research related to the project.
    \item How your project builds upon or diverges from existing work.
\end{itemize}


\section{Methodology}
\subsection{[Option 1] Model Training/Fine-tuning}
\begin{itemize}
    \item \textbf{Data Collection and Pre-processing:} Describe the dataset used, including sources, size, and any pre-processing steps undertaken (e.g., tokenization, normalization).
    \item \textbf{Model Architecture:} Specify the model used (e.g., BERT, RoBERTa, T5, GPT-2) and justify the choice based on your research objectives.
    \item \textbf{Training Parameters:} Detail the training setup, including hyperparameters like learning rate, batch size, epochs, and optimizer used.
    \item \textbf{Computational Resources:} Mention the computational resources utilized, such as GPUs, and any specific platforms (e.g., HuggingFace).
\end{itemize}

\subsection{[Option 2] Prompt Evaluation}
\begin{itemize}
    \item \textbf{Model Selection:} Identify the large language model used (e.g., GPT-4.5, Deepseek-R1) and explain the rationale for its selection.
    \item \textbf{Prompt Design:} Describe how prompts were crafted, including any iterative testing or refinement processes.
    \item \textbf{Evaluation Criteria:} Outline the criteria for evaluating the model's responses, such as accuracy, coherence, creativity, or ethical considerations.
    \item \textbf{Budget and Access:} Discuss any budgetary constraints and how access to the models was managed.
\end{itemize}

\subsection{[Option 3] Self-Proposed Projects}
\begin{itemize}
    \item \textbf{Project Description:} Provide a detailed description of your unique project idea, including objectives and expected outcomes.
    \item \textbf{Methodological Framework:} Define the framework or approach you used, including any theoretical or practical models.
    \item \textbf{Tools and Technologies:} List any specific tools, technologies, or platforms employed in your project.
    \item \textbf{Validation and Testing:} Explain how you validated your results or tested your hypotheses.
\end{itemize}

\section{Results}
\begin{itemize}
    \item Presentation of the main findings.
    \item Tables, graphs, and other visualizations to support the results.
    \item Interpretation of the results.
    \item (Optional) Comparison with previous studies. 
    \item Limitations of the study.
\end{itemize}


\section{Conclusion}
\begin{itemize}
    \item Summary of findings.
    \item Implications and potential applications.
    \item Challenges that you met.
    \item Suggestions for future research.
\end{itemize}


\section{Contributions}

For each group member, list your estimated percentage of contributions to this project (sum up to 100\%). 
\newline

\textbf{Note 1: if some group member contributes too little to the project, for example, 3\% in a group of 4, you will get a relative lower grade compared with your teammates. The student who contributes much more than all other members might have a relative higher grade compared with your teammates.}

\textbf{Note 2: the grading rubrics are given in Table 1.}

\section{Report Format}
The report should follow this template (docs/pdf) via microsoft/overleaf/latex, etc. Reference and appendix are needed. 

\bibliographystyle{plain}
\bibliography{references}
Properly formatted citations of all sources used.

\begin{table}[ht]\label{tab:rubrics}
\centering
\renewcommand{\arraystretch}{1.5}
\setlength{\tabcolsep}{4pt}

\begin{tabular}{|m{3.5cm}|m{3.5cm}|m{3.5cm}|m{3.5cm}|m{3.5cm}|}
\hline
\textbf{Project Rubrics} & \textbf{Excellent} & \textbf{Good} & \textbf{Fair} & \textbf{Poor} \\ 
\hline
\begin{tabular}[c]{@{}l@{}}Problem Definition and \\ Relevance (15\%)\end{tabular} & \textbf{(13-15 points):} Clearly defines a problem with real-world relevance that requires the LLM-related approach.  & \textbf{(10-12 points):} Defines a relevant problem with moderate complexity, integrating / leveraging LLM. & \textbf{(7-9 points):} Defines a problem, but it lacks relevance, depth, or LLM integration. & \textbf{(0-6 points):} Problem is poorly defined or irrelevant, with minimal integration of course themes. \\ 
\hline
\begin{tabular}[c]{@{}l@{}}Innovation and \\ Creativity (20\%)\end{tabular} & \textbf{(18-20 points):} Project demonstrates highly original and creative ideas, applying innovative solutions or methodologies in novel ways. & \textbf{(15-17 points):} Shows some originality and creativity, with moderate innovation in approach or solution. & \textbf{(10-14 points):} Limited creativity; applies standard methods with minimal innovation. & \textbf{(0-9 points):} Lacks creativity; relies heavily on existing solutions without any novel contributions. \\ 
\hline
\begin{tabular}[c]{@{}l@{}}Report Coherence \\ and Presentation (15\%)\end{tabular} & \textbf{(13-15 points):} Report is well-structured, with clear organization, precise language, strong coherence, and minimal errors. & \textbf{(10-12 points):} Report is well-written, with good organization and clarity, though minor issues in coherence or language may be present. & \textbf{(7-9 points):} Report conveys key points but has noticeable issues in organization, clarity, or grammar. & \textbf{(0-6 points):} Report is poorly written, with weak structure, unclear arguments, and frequent grammatical errors. \\
\hline
\begin{tabular}[c]{@{}l@{}}Technical \\ Implementation (30\%)\end{tabular} & \textbf{(26-30 points):} Implementation is robust, functional, and demonstrates advanced technical skills in the field of LLM. & \textbf{(19-25 points):} Functional implementation with moderate complexity and some advanced techniques. & \textbf{(13-18 points):} Basic implementation with limited functionality or technical depth. & \textbf{(0-12 points):} Weak or incomplete implementation with significant technical flaws. \\ 
\hline
\begin{tabular}[c]{@{}l@{}}Presentation and \\ Communication (20\%)\end{tabular} & \textbf{(17-20 points):} Clear, engaging, and well-structured presentation, effectively communicating the project's goals, methodology, and outcomes. & \textbf{(13-16 points):} Effective presentation with minor issues in clarity or structure. & \textbf{(9-12 points):} Presentation is understandable but lacks polish, clarity, or structure. & \textbf{(0-8 points):} Unclear or poorly structured presentation, failing to communicate the project effectively. \\ 
\hline
\end{tabular}
\caption{Project Rubrics}
\end{table}

\newpage


\begin{table}[ht]\label{tab:rubrics}
\centering
\renewcommand{\arraystretch}{1.5}
\setlength{\tabcolsep}{4pt}

\begin{tabular}{|m{3.5cm}|m{3.5cm}|m{3.5cm}|m{3.5cm}|m{3.5cm}|}
\hline
\textbf{Presentation Rubrics} & \textbf{Excellent} & \textbf{Good} & \textbf{Fair} & \textbf{Poor} \\ 
\hline
\begin{tabular}[c]{@{}l@{}}Presentation and \\ Communication (20\%)\end{tabular} & \textbf{(17-20 points):} Clear, engaging, and well-structured presentation, effectively communicating the project's goals, methodology, and outcomes. & \textbf{(13-16 points):} Effective presentation with minor issues in clarity or structure. & \textbf{(9-12 points):} Presentation is understandable but lacks polish, clarity, or structure. & \textbf{(0-8 points):} Unclear or poorly structured presentation, failing to communicate the project effectively. \\ 
\hline
\end{tabular}
% \caption{Project Rubrics}
\end{table}
% \begin{table}[h!]
% \centering
% \begin{tabular}{|>{\centering\arraybackslash}p{0.1\linewidth}|>{\centering\arraybackslash}p{0.3\linewidth}|>{\centering\arraybackslash}p{0.3\linewidth}|>{\centering\arraybackslash}p{0.2\linewidth}|}
% \hline
% \textbf{Order} & \textbf{Team-Name} & \textbf{Presenter} & \textbf{Grades/Scores} \\
% \hline\hline
% 1 & DeepSleep & Zihao Tang &  \\
% \hline
% 2 & BioGroup & Weicai Long &  \\
% \hline
% 3 & DriverVL & Zhijie Qiu &  \\
% \hline
% 4 & GPleaseT & Xiaorui Zhang &  \\
% \hline
% 5 & M3LM & Haowen YANG &  \\
% \hline
% 6 & DeepSick & Qi ZENG &  \\
% \hline
% 7 & Prompter & Xingjian TAO &  \\
% \hline
% 8 & LyingFlat & Jiahao ZENG &  \\
% \hline
% 9 & CodeInsightBench & Yanhong Wang &  \\
% \hline
% 10 & DeepHub & Yijie XU, Songyue Guo &  \\
% \hline
% 11 & AAAResearch & Miao Peng &  \\
% \hline
% 12 & PaperReader & Xingyu LI &  \\
% \hline
% 13 & TBD & Longhan Zhang &  \\
% \hline
% 14 & ComLLM & Ziyi Lu &  \\
% \hline
% 15 & ShallowSeek & Jingyi Pan &  \\
% \hline
% 16 & FT9 & Honghao Liu &  \\
% \hline
% 17 & SQLLLM & Huanze DONG &  \\
% \hline
% 18 & TBD & Ziyue Gao &  \\
% \hline
% 19 & JX & Liuchang JING &  \\
% \hline
% 20 & TBD & Liguo Lin &  \\
% \hline
% 21 & table LLM & Yuchen LIU &  \\
% \hline
% 22 & gymLLM & Zhongxi Tan &  \\
% \hline
% 23 & TAILab & Yuemeng Zhao &  \\
% \hline
% 24 & 404TeamFound & Jixin Yang &  \\
% \hline
% \end{tabular}
% \end{table}

\end{document}
